% This is samplepaper.tex — Springer LNCS format
\documentclass[runningheads]{llncs}

\usepackage[T1]{fontenc}
\usepackage{graphicx}
\usepackage{amsmath}
\usepackage{booktabs}
\usepackage{multirow}
\usepackage{url}

\begin{document}

%% ----------------------------------------------------------------
%%  TITLE & AUTHORS
%% ----------------------------------------------------------------
\title{[T\'ITULO POR DEFINIR]}

\author{Arturo~Cant\'u~Olivares\inst{1} \and
Diego~Sebastian~Cruz~Cervantes\inst{1} \and
Luis~Fernando~Maldonado\inst{1} \and
Jes\'us~Gabriel~Gudi\~no~Lara\inst{1} \and
Daniel Cantón Enriquez\inst{1} \and
[Colaboradores~Terraquea]\inst{2}}

\authorrunning{D.~[Apellido] et al.}

\institute{Universidad Aut\'onoma de Quer\'etaro (UAQ), Quer\'etaro, M\'exico\\
\email{dcruz54@alumnos.uaq.mx}
\and
Terraquea, M\'exico\\
\email{[correo@terraquea.com]}}

\maketitle

%% ----------------------------------------------------------------
%%  ABSTRACT
%% ----------------------------------------------------------------
\begin{abstract}
[PENDIENTE --- redactar al final cuando los resultados finales est\'en completos.]

\emph{Estructura sugerida (4--6 oraciones):}
\begin{enumerate}
  \item Contexto: la detecci\'on de \'arboles individuales (ITD) en nubes de puntos LiDAR UAV de alta densidad (ULS) es fundamental para inventarios forestales automatizados a escala.
  \item Problema: no est\'a claro si el preprocesado sem\'antico de la nube de puntos es el factor determinante en la calidad de la detecci\'on de instancias individuales.
  \item M\'etodo: comparamos tres estrategias de preprocesado --- sin preprocesado, Random Forest sobre features geom\'etricas y PointNet++ MSG --- seguidas de dos variantes de segmentador volum\'etrico Watershed~3D (seeding por CHM y seeding por densidad de copa) calibradas por flujo, evaluadas sobre FOR-instance V1 en cinco tipos de bosque.
  \item Resultado clave: el preprocesado sem\'antico mejora la precisi\'on del pipeline pero el segmentador de instancias es el cuello de botella principal; la combinaci\'on RF + Watershed~3D con seeding por densidad de copa alcanza el mejor F1 de instancia (0.209 en test, +47\% sobre RF + CHM y +87\% sobre la l\'inea base).
  \item Conclusi\'on: en datos ULS de alta densidad, Random Forest con features geom\'etricas supera a PointNet++ MSG como preprocesado sem\'antico para pipelines ITS, y el seeding por densidad de copa mejora consistentemente al seeding cl\'asico por CHM, sin requerir GPU ni arquitectura profunda.
\end{enumerate}

\keywords{Individual Tree Segmentation \and UAV laser scanning (ULS) \and Segmentaci\'on sem\'antica \and Watershed 3D \and Density seeding \and Random Forest \and PointNet++ \and FOR-instance}
\end{abstract}

%% ================================================================
%%  1. INTRODUCCI\'ON
%% ================================================================
\section{Introducci\'on}\label{sec:intro}

\subsection{Motivaci\'on}

Los inventarios forestales tradicionales dependen de mediciones de campo costosas, de baja cobertura espacial y de actualizaci\'on infrecuente. El UAV laser scanning (ULS) combina la perspectiva a\'erea del LiDAR aerotransportado con densidades de puntos uno o dos \'ordenes de magnitud superiores al ALS convencional ($\sim$7{,}000~pts/m$^2$ frente a $\sim$5~pts/m$^2$). Esta ultra-alta densidad permite caracterizar el bosque a nivel de \'arbol individual con gran detalle, incluyendo no solo copas superiores, sino tambi\'en vegetaci\'on baja, troncos, ramas y suelo. Como consecuencia, la segmentaci\'on autom\'atica de \'arboles individuales en datos ULS sigue siendo un problema abierto por el ruido potencial en modelos de altura, la ambig\"uedad entre puntos de copa y vegetaci\'on subyacente, y la variabilidad estructural entre tipos de bosque.

[PENDIENTE EXPANDIR --- 2--3 p\'arrafos adicionales sobre la importancia del problema y el estado general del campo]

\subsection{Problema espec\'ifico}

Un pipeline t\'ipico de Individual Tree Segmentation (ITS) sobre datos ULS de alta densidad comprende dos etapas en cascada: (1)~preprocesado sem\'antico, que separa los puntos de \'arbol de los de suelo y vegetaci\'on no arb\'orea, y (2)~segmentaci\'on de instancias individuales, que asigna cada punto de \'arbol a un identificador de \'arbol \'unico. En trabajos previos, la calidad del preprocesado sem\'antico suele reportarse como un factor habilitador del rendimiento de instancia, pero rara vez se a\'isla experimentalmente su contribuci\'on frente al segmentador~\cite{qi2017,breiman2001,puliti2023}.
Esta suposici\'on no ha sido evaluada sistem\'aticamente en datos ULS de alta densidad sobre m\'ultiples tipos de bosque mediante un experimento controlado donde el segmentador de instancias sea id\'entico y solo var\'ie el preprocesado.

\subsection{Contribuciones}

Este trabajo presenta las siguientes contribuciones:

\begin{enumerate}
  \item Una comparaci\'on sistem\'atica de tres estrategias de preprocesado sem\'antico --- nube cruda sin preprocesado, Random Forest sobre features geom\'etricas y PointNet++ MSG --- combinadas con dos variantes de segmentador Watershed~3D (seeding cl\'asico por CHM y seeding por densidad de copa), sobre el benchmark FOR-instance~V1~\cite{puliti2023} que cubre cinco tipos de bosque en cuatro continentes, con calibraci\'on independiente por flujo.

  \item Un an\'alisis factorial $3 \times 2$ (preprocesado $\times$ seeding) que permite descomponer la contribuci\'on de cada etapa al rendimiento del pipeline e identificar el cuello de botella dominante.

  \item Evidencia sobre la viabilidad de Random Forest con features geom\'etricas como alternativa al aprendizaje profundo para preprocesado sem\'antico en pipelines ITS sobre datos ULS de alta densidad, sin requerir GPU ni arquitectura profunda.

  \item Una variante del seeding del Watershed~3D basada en densidad de puntos en la banda superior de la copa, que mejora consistentemente la detecci\'on de semillas en canopias de altura uniforme (bosques boreales densos) frente al seeding cl\'asico por CHM.
\end{enumerate}

\subsection{Estructura del paper}

El resto del paper se organiza de la siguiente manera. La Secci\'on~\ref{sec:related} revisa el trabajo relacionado. La Secci\'on~\ref{sec:dataset} describe el dataset utilizado. La Secci\'on~\ref{sec:method} detalla los cinco pipelines propuestos (tres preprocesados $\times$ dos variantes de segmentador) y el detector compartido. La Secci\'on~\ref{sec:eval} presenta el protocolo de evaluaci\'on. La Secci\'on~\ref{sec:results} reporta y analiza los resultados. La Secci\'on~\ref{sec:discussion} discute las implicaciones y limitaciones. La Secci\'on~\ref{sec:conclusions} concluye.

%% ================================================================
%%  2. TRABAJO RELACIONADO
%% ================================================================
\section{Trabajo Relacionado}\label{sec:related}

\subsection{Detecci\'on de \'arboles individuales en LiDAR aerotransportado}

% [PENDIENTE]
% Puntos a cubrir:
% - CHM-watershed como m\'etodo cl\'asico: Silva et al. 2016, Dalponte & Coomes 2016, Popescu et al. 2002
% - CHM-watershed en ALS convencional (baja densidad)
% - Diferencia de r\'egimen entre ALS convencional y ULS denso
% - Benchmarks ALS/ULS: NeonTreeEvaluation y FOR-instance
% - Transici\'on natural hacia m\'etodos de LiDAR a\'ereo denso (ULS / HD-ALS)

\subsection{Detecci\'on de \'arboles individuales en LiDAR denso (ULS / HD-ALS)}

% [PENDIENTE]
% Puntos a cubrir:
% - M\'etodos ITD sobre ULS/HD-ALS y CHM de alta resoluci\'on
% - Efecto de la ultra-alta densidad en detecci\'on por m\'aximos locales
% - Compromiso entre suavizado, resoluci\'on y fusi\'on de copas
% - Gap: pocos trabajos a\'islan el impacto del preprocesado sem\'antico frente al detector
% - Relaci\'on entre calidad sem\'antica y ITD en ULS de alta densidad

\subsection{Clasificaci\'on sem\'antica de nubes de puntos forestales}

% [PENDIENTE]
% Puntos a cubrir:
% - M\'etodos cl\'asicos con features manuales: Random Forest, SVM
% - M\'etodos de aprendizaje profundo: PointNet \cite{qi2016}, PointNet++ \cite{qi2017}
% - Aplicaciones de PointNet++ en datos forestales ULS / airborne densos
% - Posicionamiento: la mayor\'ia eval\'ua la clasificaci\'on como tarea final, no como etapa intermedia de un pipeline ITD en datos ULS/airborne densos

\subsection{Benchmarks de segmentaci\'on forestal}

% [PENDIENTE]
% Puntos a cubrir:
% - FOR-instance V1 \cite{puliti2023}
% - FOR-instanceV2 \cite{xiang2025}
% - NeonTreeEvaluation \cite{weinstein2021}
% - Posicionamiento de FOR-instance como opci\'on natural para evaluar pipelines ULS end-to-end

%% ================================================================
%%  3. DATASET
%% ================================================================
\section{Dataset}\label{sec:dataset}

\subsection{FOR-instance}

El dataset FOR-instance~\cite{puliti2023} es un benchmark de acceso abierto dise\~nado para la evaluaci\'on de m\'etodos de segmentaci\'on sem\'antica e de instancias de \'arboles individuales en nubes de puntos LiDAR de alta densidad. La versi\'on V1 comprende seis colecciones (NIBIO, NIBIO2, CULS, TUWIEN, RMIT, SCION) de datos adquiridos con esc\'aneres l\'aser Riegl (VUX-1 UAV y Mini-VUX) montados en veh\'iculos a\'ereos no tripulados (UAV laser scanning, ULS) sobre sitios geogr\'aficamente diversos que cubren tipos de bosque representativos a nivel global.

En este trabajo se utilizan las cinco colecciones que cubren tipos de bosque distintos (NIBIO, CULS, TUWIEN, RMIT, SCION). Se excluye expl\'icitamente NIBIO2 por una raz\'on metodol\'ogica: NIBIO2 corresponde a parcelas adicionales del mismo tipo de bosque boreal con\'ifero noruego (\emph{Picea abies}) muestreado por NIBIO. Incluir ambas colecciones sobre-representar\'ia el bosque boreal en el conjunto de evaluaci\'on, sesgando las m\'etricas agregadas hacia un \'unico tipo estructural y diluyendo la se\~nal experimental sobre la generalizaci\'on entre tipos de bosque, que es objetivo central del estudio.

Las nubes de puntos presentan densidades que var\'ian fuertemente entre colecciones: medidas sobre los 32 plots empleados en este trabajo, el rango global es $454$--$10\,156$~pts/m$^2$ (mediana por colecci\'on: RMIT $473$, TUWIEN $1\,406$, CULS $2\,510$, SCION $3\,064$, NIBIO $7\,246$), todas ellas considerablemente m\'as altas que el LiDAR aerotransportado convencional (${\sim}5$~pts/m$^2$) y consistentes con los rangos reportados por Puliti et al.~\cite{puliti2023}. Esta alta densidad permite la anotaci\'on manual precisa de instancias individuales. Cada punto est\'a anotado con una etiqueta sem\'antica de entre seis clases --- suelo, vegetaci\'on baja, tronco, ramas vivas, ramas muertas y puntos de borde (outpoints) --- as\'i como con un identificador de \'arbol individual (treeID) para todos los \'arboles del dosel anotados.

\subsection{Colecciones y tipos de bosque}

El dataset agrupa los plots en cinco colecciones geogr\'aficas, cada una representando un tipo de bosque distinto (Tabla~\ref{tab:collections}).

\begin{table}[t]
\centering
\caption{Colecciones del dataset FOR-instance V1 utilizadas en este trabajo.}\label{tab:collections}
\begin{tabular}{llccc}
\toprule
C\'odigo & Pa\'is & Tipo de bosque & Plots test & \'Arboles GT \\
\midrule
NIBIO  & Noruega       & Boreal con\'ifero (\emph{P.~abies})        & 6 & 161 \\
CULS   & Rep.\ Checa   & Templado con\'ifero (\emph{P.~sylvestris}) & 1 &  20 \\
TUWIEN & Austria       & Mixto deciduo aluvial                      & 1 &  35 \\
RMIT   & Australia     & Escler\'ofilo seco (eucaliptos)             & 1 &  64 \\
SCION  & Nueva Zelanda & Plantaci\'on con\'ifera (\emph{P.~radiata}) & 2 &  43 \\
\midrule
\textbf{Total} & & & \textbf{11} & \textbf{323} \\
\bottomrule
\end{tabular}
\end{table}

La diversidad geogr\'afica y estructural de las colecciones es relevante para evaluar la generalizaci\'on de los m\'etodos: NIBIO representa el 50\% del conjunto de test y corresponde a bosque boreal de alta densidad, mientras que RMIT y TUWIEN representan estructuras de copa radicalmente diferentes entre s\'i y respecto al bosque boreal.

\subsection{Splits y desdoblamiento interno}\label{sec:splits}

El dataset prescribe un split oficial: 70\% de los plots para desarrollo (dev) y 30\% para evaluaci\'on (test), con todos los sitios geogr\'aficos representados en ambas particiones. En este trabajo se utiliza el split oficial sin modificaciones. El conjunto dev se emplea para entrenamiento de modelos (Flujos~B y~C) y calibraci\'on de hiperpar\'ametros del detector (los tres flujos). El conjunto test se utiliza exclusivamente para reportar las m\'etricas finales y no interviene en ninguna decisi\'on de dise\~no o calibraci\'on.

\paragraph{Desdoblamiento interno dev~$\rightarrow$~train+val.}
Dado que FOR-instance~V1 define exclusivamente las particiones dev y test en su metadata oficial, el conjunto dev se subdivide internamente en un conjunto de entrenamiento (80\% de los plots) y un conjunto de validaci\'on interno (20\% restante) mediante una partici\'on aleatoria a nivel de plot con semilla fija (\texttt{random\_state}~$= 42$) para reproducibilidad. La partici\'on se realiza a nivel de archivo \texttt{.las} y no a nivel de punto para evitar \emph{data leakage} espacial entre train y val, siguiendo el criterio adoptado por ForAINet~\cite{henrich2024} y SegmentAnyTree~\cite{wielgosz2024}. El conjunto val interno se utiliza para: (i) monitorear la p\'erdida durante el entrenamiento de PointNet++ y guiar el scheduler, (ii) ejecutar el grid search de los hiperpar\'ametros del Watershed~3D de forma independiente por flujo, y (iii) reportar las m\'etricas val que se presentan en la Secci\'on~\ref{sec:results}. El conjunto test oficial de FOR-instance permanece intacto y se reserva exclusivamente para el reporte de m\'etricas finales, sin intervenir en ninguna decisi\'on de dise\~no o calibraci\'on.

\paragraph{Reproducibilidad.}
Para garantizar reproducibilidad, todas las fuentes de aleatoriedad del pipeline --- partici\'on train/val interna, submuestreo por plot del Random Forest, inicializaci\'on y barajado del entrenamiento de PointNet++, e inferencia multi-pasada de PointNet++ --- usan la semilla fija $42$ salvo indicaci\'on expl\'icita en contrario. En el caso del submuestreo por plot del Random Forest, la semilla efectiva se deriva de forma determin\'istica del nombre del plot v\'ia \texttt{CRC32}, de modo que cada plot reciba siempre el mismo subconjunto de puntos con independencia del orden de iteraci\'on.

\subsection{Ground truth para segmentaci\'on de instancias}

Para la tarea de segmentaci\'on de \'arboles individuales el ground truth est\'a definido a nivel de punto: cada punto de la nube tiene asociado un identificador entero \texttt{treeID} que indica a qu\'e \'arbol individual pertenece (o cero para puntos no arb\'oreos). El protocolo de evaluaci\'on compara el conjunto de puntos predicho para cada instancia contra el conjunto de puntos GT del \texttt{treeID} mediante IoU 3D punto a punto (Secci\'on~\ref{sec:eval}). Los puntos clasificados como \emph{outpoint} en FOR-instance --- correspondientes a \'arboles parcialmente observados en los bordes de los plots --- se excluyen de la evaluaci\'on tanto para predicciones como para ground truth.

%% ================================================================
%%  4. METODOLOG\'IA
%% ================================================================
\section{Metodolog\'ia}\label{sec:method}

\subsection{Visi\'on general del pipeline}

Los flujos propuestos comparten una arquitectura de dos etapas en cascada. La Etapa~1 es el preprocesado sem\'antico, cuyo objetivo es separar los puntos de \'arbol del resto de la nube (suelo, vegetaci\'on baja, ramas ca\'idas). La Etapa~2 es el segmentador de instancias, que opera sobre los puntos clasificados como \'arbol y asigna a cada uno un identificador de \'arbol individual. El dise\~no experimental var\'ia ambas etapas de forma factorial: tres estrategias de preprocesado sem\'antico (Etapa~1) combinadas con dos variantes del segmentador Watershed~3D (Etapa~2), produciendo cinco flujos evaluados. Los par\'ametros del segmentador se calibran mediante grid search independiente por flujo sobre el split val maximizando el F1 de instancia (umbral IoU 3D~$\geq 0.5$).

Esta calibraci\'on independiente garantiza que cada flujo opere bajo sus condiciones \'optimas, de modo que las diferencias en F1 de instancia entre flujos reflejen el efecto del preprocesado sem\'antico y de la estrategia de seeding, y no diferencias en la configuraci\'on del segmentador.

Las tres estrategias de preprocesado sem\'antico (Etapa~1) son:
\begin{itemize}
  \item \textbf{Sin preprocesado (baseline):} nube cruda completa.
  \item \textbf{Random Forest:} clasificador sobre features geom\'etricas (ML cl\'asico).
  \item \textbf{PointNet++ MSG:} clasificador de aprendizaje profundo.
\end{itemize}

Las dos variantes del segmentador (Etapa~2) son:
\begin{itemize}
  \item \textbf{Watershed~3D--CHM:} seeding sobre la envolvente superior del volumen (equivalente a un CHM local).
  \item \textbf{Watershed~3D--Density:} seeding sobre una superficie de densidad de puntos en la banda superior de la copa.
\end{itemize}

Los cinco flujos evaluados son:
\begin{itemize}
  \item \textbf{Flujo~A:} Sin preprocesado + Watershed~3D--CHM (baseline).
  \item \textbf{Flujo~B:} Random Forest + Watershed~3D--CHM.
  \item \textbf{Flujo~C:} PointNet++ MSG + Watershed~3D--CHM.
  \item \textbf{Flujo~D:} Random Forest + Watershed~3D--Density.
  \item \textbf{Flujo~E:} PointNet++ MSG + Watershed~3D--Density.
\end{itemize}

Los flujos D y E reutilizan las predicciones sem\'anticas de los flujos B y C respectivamente; la \'unica diferencia es la variante del segmentador de instancias. Todos los flujos operan sobre la nube completa del plot sin dividirla en bloques: el clasificador se eval\'ua punto a punto sobre la nube entera y el segmentador opera sobre el subconjunto de puntos resultante (todos los puntos v\'alidos para el Flujo~A; los clasificados como \'arbol para los dem\'as).

\subsection{Flujo~A --- Sin preprocesado (baseline)}

El Flujo~A no aplica ning\'un clasificador sem\'antico. La nube de puntos completa del plot, conservando todos los puntos con etiqueta de clasificaci\'on de FOR-instance en el conjunto $\{1, 2, 4, 5, 6\}$ (suelo y vegetaci\'on baja no arb\'orea, y todas las variantes de \'arbol), se pasa directamente al segmentador Watershed~3D de la Etapa~2. Los puntos con clase $\in \{0, 3\}$ (\emph{never-classified} y \emph{outpoint}, estos \'ultimos correspondientes a \'arboles parcialmente observados en los bordes del plot) se excluyen en todos los flujos por protocolo, tanto en el preprocesado como en la evaluaci\'on de instancias, siguiendo la especificaci\'on oficial de FOR-instance. Este flujo establece el nivel de rendimiento de instancia alcanzable sin ninguna separaci\'on sem\'antica previa y sirve como referencia para cuantificar el valor a\~nadido de los clasificadores de los Flujos~B y~C.

\paragraph{Preprocesado com\'un de altura sobre el suelo (HAG).}
Aunque el Flujo~A no tiene un clasificador sem\'antico, comparte con los Flujos~B y~C el mismo preprocesado geom\'etrico de altura sobre el suelo, necesario para que el segmentador Watershed~3D opere en coordenadas invariantes al relieve (Secci\'on~\ref{sec:watershed}). El DTM utilizado para calcular la HAG se construye por plot tomando la mediana de $z$ por celda sobre los puntos clasificados como suelo (clase~2 de FOR-instance), rasterizado a una resoluci\'on horizontal de $0{,}5$\,m. Las celdas sin puntos de suelo se rellenan por interpolaci\'on bilineal (\texttt{scipy.interpolate.griddata}, m\'etodo \texttt{linear}), con \emph{fallback} al vecino m\'as cercano para las celdas de borde que queden fuera del casco convexo de los puntos de suelo. Posteriormente se aplica un suavizado gaussiano cuyo $\sigma$ equivale a una ventana $5 \times 5$ ($\sigma \approx 2{,}5$ celdas) para reducir artefactos locales. Tras calcular la HAG, el canal se trunca al rango $[0, h_{\max}]$ con $h_{\max} = 50$\,m: los valores negativos (ruido del DTM bajo el suelo) se clampean a cero y los valores superiores se saturan a $50$\,m, un umbral que cubre con holgura las alturas m\'aximas observadas en los cinco sitios de FOR-instance~V1. Para el Flujo~C (PointNet++) este canal HAG se renormaliza adicionalmente al intervalo $[0, 1]$ dividiendo por $h_{\max}$, de modo que sea comparable con la intensidad l\'aser (tambi\'en en $[0, 1]$) como feature de entrada a la red.

\subsection{Flujo~B --- Random Forest sobre features geom\'etricas}

\subsubsection{Features.}
Cada punto se representa mediante un vector de 27 features geom\'etricas: 13 descriptores dependientes de escala calculados a dos vecindades ($k\!=\!20$ y $k\!=\!50$ vecinos), m\'as una feature adicional de densidad volum\'etrica compartida entre escalas ($13 \times 2 + 1 = 27$). La selecci\'on sigue el cat\'alogo can\'onico de Weinmann et al.~\cite{weinmann2017}, complementado con descriptores verticales basados en la altura sobre el suelo (HAG) que han demostrado ser predictores dominantes de la separaci\'on \'arbol/no-\'arbol en bosques densos~\cite{li2023,bremer2023}.

Por cada vecindad de tama\~no $k$ se calculan ocho descriptores eigen-basados a partir de los autovalores normalizados $\lambda_1 \geq \lambda_2 \geq \lambda_3$ de la matriz de covarianza local (estimada con divisor $k$, no $k-1$, siguiendo la convenci\'on de Weinmann et al.~\cite{weinmann2015}) --- linealidad, planaridad, esfericidad, omnivarianza, anisotrop\'ia, eigenentrop\'ia, suma de eigenvalores y cambio de curvatura --- m\'as cinco descriptores locales: HAG del punto central, verticalidad (definida como $1 - |\langle \mathbf{n}, [0,0,1]\rangle|$ con $\mathbf{n}$ la normal local), rugosidad ($\text{std}(z)$ entre vecinos), rango vertical ($\max(z)-\min(z)$ entre vecinos) e intensidad normalizada del retorno l\'aser. La feature restante --- densidad volum\'etrica local --- es independiente de la escala $k$ y se calcula una sola vez por punto como la cardinalidad de la bola de radio fijo $r$ centrada en el punto. Las dos escalas $k$ proporcionan contexto fino y grueso simult\'aneamente, siguiendo la estrategia multi-escala de Weinmann et al.~\cite{weinmann2015}.

Para evitar que en plots de baja densidad (p.\,ej.\ RMIT, mediana $\sim\!473$\,pts/m$^2$) los vecinos queden demasiado alejados y contaminen los descriptores locales con contexto irrelevante, se impone un radio m\'aximo de b\'usqueda de $5{,}0$\,m: los vecinos a mayor distancia dentro de los $k$ m\'as pr\'oximos se descartan, y los puntos con menos de cuatro vecinos v\'alidos tras el corte reciben \emph{features} en cero para esa escala (valor neutro para el Random Forest). Este r\'egimen afecta a una fracci\'on muy reducida de puntos en los plots considerados y se reporta expl\'icitamente para trazabilidad.

\begin{table}[t]
\centering
\caption{Cat\'alogo de features geom\'etricas para el clasificador Random Forest (Flujo~B). 13 descriptores dependientes de escala calculados a dos vecindades $k\!=\!20$ y $k\!=\!50$, m\'as 1 feature de densidad volum\'etrica compartida entre escalas, producen un vector total de $13 \times 2 + 1 = 27$ dimensiones por punto. En las filas 1--6 y 8 los $\lambda_i$ denotan los autovalores normalizados $\lambda_i / (\lambda_1+\lambda_2+\lambda_3)$, siguiendo la convenci\'on de Weinmann et al.~\cite{weinmann2015}; la fila~7 (suma de eigenvalores) usa los autovalores crudos.}\label{tab:rf-features}
\begin{tabular}{lll}
\toprule
\# & Feature & F\'ormula / definici\'on \\
\midrule
\multicolumn{3}{l}{\emph{Por escala $k \in \{20, 50\}$}} \\
1  & Linealidad             & $(\lambda_1 - \lambda_2)/\lambda_1$       \\
2  & Planaridad             & $(\lambda_2 - \lambda_3)/\lambda_1$       \\
3  & Esfericidad            & $\lambda_3/\lambda_1$                     \\
4  & Omnivarianza           & $(\lambda_1\lambda_2\lambda_3)^{1/3}$     \\
5  & Anisotrop\'ia          & $(\lambda_1 - \lambda_3)/\lambda_1$       \\
6  & Eigenentrop\'ia        & $-\sum_i (\lambda_i/\!\sum_j\!\lambda_j)\,\ln(\lambda_i/\!\sum_j\!\lambda_j)$ \\
7  & Suma de eigenvalores   & $\lambda_1 + \lambda_2 + \lambda_3$       \\
8  & Cambio de curvatura    & $\lambda_3/(\lambda_1+\lambda_2+\lambda_3)$ \\
9  & HAG                    & $z - \text{DTM}(x,y)$                     \\
10 & Verticalidad           & $1 - |\langle \mathbf{n}, [0,0,1]\rangle|$ \\
11 & Rugosidad              & $\text{std}(z_{\text{vecinos}})$          \\
12 & Rango vertical         & $\max(z)-\min(z)$ entre vecinos           \\
13 & Intensidad normalizada & $I / I_{\max}$ por plot                   \\
\midrule
\multicolumn{3}{l}{\emph{Compartida entre escalas}} \\
14 & Densidad volum\'etrica local & $|\{p_j : \lVert p_j - p_i \rVert \leq r\}|$ (cardinalidad de la bola de radio $r$) \\
\bottomrule
\end{tabular}
\end{table}

\paragraph{Nota sobre la feature de densidad volum\'etrica.}
La feature de densidad volum\'etrica (fila~14) se computa como la cardinalidad real del conjunto $\{p_j : \lVert p_j - p_i \rVert \leq r\}$ con $r = \texttt{density\_radius} = 0{,}5$\,m, id\'entico para todos los plots. A diferencia de los descriptores eigen-basados, no depende de la vecindad $k$: se calcula una sola vez por punto mediante una consulta de bola al \texttt{KDTree} global del plot (\texttt{cKDTree.query\_ball\_point}), de modo que su valor no satura por arriba en zonas de alta densidad. Se reporta en forma de conteo entero para coincidir con la convenci\'on de \emph{neighborhood cardinality} de Weinmann et al.~\cite{weinmann2015}; al ser una transformaci\'on lineal de la densidad volum\'etrica $|\cdot|/V_{\text{esfera}}$, el factor $1/V_{\text{esfera}}$ es absorbido por los umbrales de split del Random Forest sin afectar las decisiones.

\subsubsection{Hiperpar\'ametros y entrenamiento.}
El clasificador es un \texttt{RandomForestClassifier} de scikit-learn entrenado a nivel de plot (no por bloques): los 27 features se computan sobre la nube completa de cada plot del split train, los puntos se etiquetan como \'arbol/no-\'arbol seg\'un el ground truth de FOR-instance, y se aplica un submuestreo balanceado \emph{por plot} que retiene hasta $12\,000$ puntos por clase en cada plot. Este esquema equipara la contribuci\'on de cada sitio al conjunto final de entrenamiento, evitando que los sitios con plots grandes (NIBIO) dominen el modelo. La alternativa de submuestrear globalmente tras concatenar todos los plots introducir\'ia un sesgo proporcional al n\'umero de puntos por sitio, contrario al objetivo de evaluar generalizaci\'on entre tipos de bosque que motiva el presente trabajo. El submuestreo es determin\'istico: la semilla efectiva de cada plot se deriva de forma determin\'istica del nombre del archivo v\'ia \texttt{CRC32} sobre la semilla global \texttt{random\_state}$=42$, de forma que cada plot siempre recibe los mismos puntos con independencia del orden de iteraci\'on. El modelo se entrena con $\texttt{n\_estimators}=200$ \'arboles, sin l\'imite de profundidad ($\texttt{max\_depth}=\texttt{None}$) y con $\texttt{class\_weight}=\texttt{balanced}$ para compensar el desbalance residual.

La elecci\'on $\texttt{max\_depth}=\texttt{None}$ corresponde al ajuste por defecto de Random Forest seg\'un Breiman~\cite{breiman2001}: cada \'arbol crece hasta la pureza de la hoja, y el control de varianza recae sobre el promediado del bagging y la aleatorizaci\'on de features ($\texttt{max\_features}=\sqrt{p}$). Limitar la profundidad introducir\'ia sesgo sin un beneficio claro de generalizaci\'on para la dimensionalidad y el tama\~no de muestra empleados. La elecci\'on $\texttt{n\_estimators}=200$ se valid\'o emp\'iricamente como punto de estabilizaci\'on de la curva Out-Of-Bag.

\begin{table}[t]
\centering
\caption{Hiperpar\'ametros del clasificador Random Forest (Flujo~B).}\label{tab:rf-hp}
\begin{tabular}{lcc}
\toprule
Hiperpar\'ametro & Valor & Criterio de selecci\'on \\
\midrule
n\_estimators    & 200       & Estabilizaci\'on curva OOB       \\
max\_depth       & None      & Default Breiman~\cite{breiman2001} \\
max\_features    & sqrt      & Reducci\'on de varianza~\cite{breiman2001} \\
class\_weight    & balanced  & Desbalance \'arbol/no-\'arbol      \\
max\_samples\_per\_plot & 12{,}000  & Submuestreo por plot, equilibrio clases y sitios \\
\bottomrule
\end{tabular}
\end{table}

\paragraph{Inferencia.} En inferencia el clasificador se aplica directamente a la nube completa del plot: los 27 features se computan punto a punto y la predicci\'on binaria se obtiene aplicando el modelo entrenado, sin agregaci\'on por bloques ni promediado de probabilidades.

\subsection{Flujo~C --- PointNet++ MSG}

\subsubsection{Arquitectura.}
Se emplea la arquitectura PointNet++~\cite{qi2017} en su variante \emph{Multi-Scale Grouping} (MSG) para segmentaci\'on sem\'antica punto a punto, con cuatro m\'odulos de agrupamiento jer\'arquico (Set Abstraction) seguidos de cuatro m\'odulos de propagaci\'on de features (Feature Propagation). La implementaci\'on toma como base el repositorio p\'ublico de Xu~Yan~\cite{yan2019pointnet2pytorch}. Los radios de b\'usqueda de las capas \emph{Set Abstraction} operan en el \emph{espacio normalizado} del sample (tras dividir las coordenadas por el radio m\'aximo del bloque de 8{,}192 puntos), no en metros absolutos; esto hace que la escala efectiva de cada radio dependa del tama\~no f\'isico del plot. La red recibe como entrada nubes submuestreadas a 8{,}192 puntos por plot, donde cada punto se representa por un vector de cinco dimensiones: las tres coordenadas $XYZ$ normalizadas, la altura sobre el suelo HAG normalizada y la intensidad del retorno l\'aser. La cabeza de clasificaci\'on punto a punto consiste en una capa lineal $128 \to 64$ con \texttt{BatchNorm}, \texttt{ReLU} y \emph{dropout} con $p = 0{,}5$, seguida de una capa lineal $64 \to C$ con $C = 2$ clases. La salida son log-probabilidades por punto sobre las dos clases (\'arbol / no-\'arbol).

\paragraph{Justificaci\'on de la entrada de 5 dimensiones.}
La elecci\'on de $XYZ + \text{HAG} + I$ como entrada se basa en mantener \emph{consistencia con las se\~nales crudas disponibles para los otros flujos}. El Flujo~B (RF) consume HAG e intensidad como features de primer orden; usar las mismas se\~nales en el Flujo~C garantiza que la comparaci\'on entre RF y PointNet++ aisle el efecto del modelo (features manuales vs.\ aprendidas) y no el efecto de la informaci\'on disponible en la entrada. Esta es la \'unica forma honesta de comparar dos clasificadores sem\'anticos en condiciones equivalentes: ambos ven los mismos canales crudos. Esta elecci\'on es adem\'as consistente con Wang et al.~\cite{wang2023}, que reporta mejoras significativas al incorporar HAG e intensidad en PointNet++ para clasificaci\'on forestal.

\begin{table}[t]
\centering
\caption{Configuraci\'on de PointNet++ MSG (Flujo~C).}\label{tab:pointnet}
\begin{tabular}{lc}
\toprule
Par\'ametro & Valor \\
\midrule
Variante                  & MSG (Multi-Scale Grouping) \\
M\'odulos Set Abstraction & 4 \\
M\'odulos Feature Prop.\  & 4 \\
Puntos de entrada         & 8{,}192 por plot \\
Dimensi\'on de entrada    & 5 ($XYZ$ + HAG\_norm + intensidad) \\
N\'umero de clases        & 2 (\'arbol / no-\'arbol) \\
Funci\'on de p\'erdida    & NLLLoss ponderada por clase \\
Optimizador               & Adam, lr\,=\,0.001, weight\_decay\,=\,$10^{-4}$ \\
Scheduler                 & StepLR, $\gamma\,=\,0.7$ cada 20 \'epocas \\
Gradient clipping         & $\lVert g \rVert_2 \leq 1{,}0$ (norma global) \\
Dropout (cabeza FP)       & $p = 0{,}5$ \\
Batch size                & 16 \\
\'Epocas                  & 100 \\
Mixed precision           & S\'i (\texttt{torch.amp}, \texttt{bfloat16}) \\
Hardware                  & Apple M4 Pro (memoria unificada) \\
\bottomrule
\end{tabular}
\end{table}

\subsubsection{Entrenamiento.}
La p\'erdida es \texttt{NLLLoss} con pesos $w_c = N / (2 N_c + \epsilon)$, donde $N$ es el total de puntos de entrenamiento, $N_c$ el conteo de la clase $c$ y $\epsilon = 10^{-6}$, para compensar el desbalance arb\'orea/no-arb\'orea. Se optimiza con Adam (lr inicial $0{,}001$, weight decay $10^{-4}$). El scheduler reduce la tasa de aprendizaje un factor $\gamma = 0{,}7$ cada 20 \'epocas, hasta 100 \'epocas en total. Para estabilizar el entrenamiento bajo \emph{mixed precision} se aplica \emph{gradient clipping} global por norma $\ell_2$ con umbral $1{,}0$ (\texttt{torch.nn.utils.clip\_grad\_norm\_}).

La augmentaci\'on incluye tres operaciones aplicadas por muestra tras la normalizaci\'on per-sample del bloque de 8{,}192 puntos (centrado en el centroide y divisi\'on por el radio m\'aximo): (i) \emph{jitter} gaussiano independiente por coordenada con $\sigma = 0{,}02$ en el \emph{espacio normalizado} --- equivalente a una fracci\'on peque\~na del radio m\'aximo del sample y no a una distancia fija en metros ---, (ii) reflexi\'on aleatoria sobre el eje $X$ con probabilidad $0{,}5$, aprovechando la simetr\'ia horizontal aproximada de las copas forestales, y (iii) escalado isotr\'opico uniforme en $[0{,}9,\, 1{,}1]$. No se aplica rotaci\'on en $Z$ porque la orientaci\'on del grid del dataset es consistente entre plots. El batch size se fija a 16 aprovechando la memoria unificada del Apple M4 Pro, y se emplea \emph{mixed precision} en \texttt{bfloat16} (sin \emph{loss scaling}, ya que \texttt{bfloat16} preserva el rango de exponente de \texttt{float32}) para acelerar el entrenamiento.

\paragraph{Decisi\'on de entrenamiento \'unico.}
PointNet++ se entrena con una \'unica configuraci\'on de hiperpar\'ametros, sin b\'usqueda exhaustiva. El objetivo del paper no es encontrar la configuraci\'on \'optima de PointNet++ sino evaluar el valor del aprendizaje profundo como estrategia de preprocesado en condiciones representativas de uso real. En aplicaciones forestales aplicadas, un usuario que adopta PointNet++ entrena el modelo con una configuraci\'on est\'andar documentada en la literatura, no mediante optimizaci\'on intensiva. Evaluar PointNet++ bajo esas condiciones es m\'as informativo para la comunidad que reportar su mejor resultado te\'orico posible.

\paragraph{Inferencia.}
PointNet++ fue entrenado con submuestreo aleatorio global del plot completo a $N_s = 8\,192$ puntos por muestra y normalizaci\'on per-sample por centroide y radio m\'aximo. Previo a la normalizaci\'on per-sample, las coordenadas $XY$ se centran en el centroide del plot para reducir el error num\'erico de \texttt{float32} sobre coordenadas UTM (esta traslaci\'on global no altera el resultado de la normalizaci\'on per-sample posterior, que absorbe cualquier desplazamiento r\'igido). La inferencia replica exactamente esta estrategia para garantizar consistencia entre train y test: cada plot se procesa mediante $n_{\text{passes}}$ submuestreos aleatorios independientes, con
\begin{equation}
  n_{\text{passes}} \;=\; \min\!\bigl(1000,\ \max\!\bigl(50,\ \lfloor 3 N / N_s \rfloor\bigr)\bigr),
\end{equation}
donde $N$ es el n\'umero de puntos v\'alidos del plot. En cada pasada, los $N_s$ puntos seleccionados se pre-centran y normalizan como en entrenamiento, y las log-probabilidades por punto se acumulan. Tras todas las pasadas, cada punto recibe la media de las probabilidades predichas en las pasadas en que fue seleccionado, y la clase final se obtiene como $\arg\max$. Con este esquema la cobertura efectiva supera el 99\% de los puntos del plot en todos los casos observados. Los puntos residuales no cubiertos se asignan por interpolaci\'on KNN ponderada por el inverso de la distancia ($w_j \propto 1/d_j$) con $k = 3$ desde los puntos ya cubiertos m\'as cercanos.

La alternativa de una \'unica pasada con extrapolaci\'on 1-NN fue descartada por dos razones: (a)~introducir\'ia un \emph{mismatch de escala} entre train y test, ya que los radios de las capas \emph{Set Abstraction} est\'an calibrados a la escala del sample normalizado de $N_s = 8\,192$ puntos tras la divisi\'on por el radio m\'aximo del bloque, y (b)~en plots t\'ipicos etiquetar\'ia alrededor del $99\%$ de los puntos por vecindad y no por predicci\'on directa de la red, degradando la calidad sem\'antica sin justificaci\'on.

\subsection{Segmentador de instancias --- Watershed 3D volum\'etrico}\label{sec:watershed}

\subsubsection{Justificaci\'on del m\'etodo.}
La densidad ultra-alta de los datos ULS ($\sim$7{,}000~pts/m$^2$) permite ejecutar el paso de \emph{fill} del watershed directamente sobre una grilla 3D de densidad, preservando la estructura vertical de la nube. Un watershed 2D tradicional sobre un Canopy Height Model (CHM) asigna cada punto en funci\'on de su proyecci\'on $XY$, lo que produce errores sistem\'aticos cuando puntos de tronco o de rama baja caen bajo la proyecci\'on horizontal de una copa vecina. En este trabajo se emplea en cambio un \emph{Watershed~3D volum\'etrico} en el que las semillas se detectan sobre la \emph{envolvente superior del volumen} --- equivalente a un CHM derivado localmente del propio grid 3D de densidad --- pero el \emph{fill} del watershed se ejecuta sobre la grilla 3D con una m\'ascara de v\'oxeles ocupados, de modo que cada punto es asignado al \'arbol cuyo v\'oxel tridimensional lo contiene y no seg\'un su mera proyecci\'on horizontal. Esta distinci\'on es cr\'itica: el verdadero aporte del enfoque 3D no es el seeding (los m\'aximos locales de densidad 3D caer\'ian en los troncos y ramas gruesas, no en los \'apices, arruinando la segmentaci\'on; por eso la literatura ITS est\'andar usa envolventes tipo CHM para seeding incluso con fill 3D~\cite{dalponte2016}), sino el fill volum\'etrico, que es lo que realmente separa troncos adyacentes y puntos de sotobosque.

\paragraph{Uso de HAG como eje vertical.}
Una caracter\'istica clave del segmentador es que reemplaza la coordenada $Z$ absoluta por la altura sobre el suelo (Height Above Ground, HAG) antes de voxelizar la nube. La HAG se calcula sustrayendo a cada punto la altura del DTM extra\'ido del propio plot mediante filtrado de puntos clasificados como suelo. Operar en HAG produce \emph{invarianza al relieve}: dos \'arboles vecinos sobre una ladera mantienen la misma altura efectiva en HAG aunque sus picos en $Z$ absoluto difieran varios metros, evitando que la voxelizaci\'on los separe artificialmente. Es la pr\'actica est\'andar en segmentaci\'on volum\'etrica forestal en terreno con relieve no trivial.

\subsubsection{Procedimiento.}
El segmentador se aplica sobre el subconjunto de puntos clasificados como \'arbol en la Etapa~1 (o sobre la nube completa de puntos v\'alidos en el Flujo~A) y sigue cuatro pasos:

\paragraph{Paso~1 --- Voxelizaci\'on en $(X, Y, \text{HAG})$.} Los puntos seleccionados se voxelizan en una grilla 3D regular en la que el eje vertical es HAG en lugar de $Z$ absoluto. El tama\~no de v\'oxel es par\'ametro del grid search.

\paragraph{Paso~2 --- Suavizado gaussiano 3D.} Sobre la grilla de densidad por v\'oxel se aplica un filtro gaussiano 3D para eliminar fluctuaciones de alta frecuencia debidas a ramas individuales y ruido de medici\'on. El par\'ametro $\sigma$ es parte del grid search.

\paragraph{Paso~3 --- Detecci\'on de semillas.} Este paso admite dos variantes:

\emph{Variante CHM (Flujos A, B, C).} Para cada columna $(x, y)$ del grid 3D se identifica el v\'oxel ocupado m\'as alto, construyendo un mapa de alturas local (equivalente a un CHM derivado del propio grid). Sobre este mapa se aplica \texttt{peak\_local\_max} con separaci\'on m\'inima entre picos dada por \emph{min\_crown\_radius\_m} / \emph{voxel\_size} y umbral absoluto igual a \emph{min\_tree\_height}. Las semillas con altura efectiva menor a \emph{min\_tree\_height} se descartan.

\emph{Variante Density (Flujos D, E).} En bosques densos con canopias de altura uniforme (e.g.\ bosque boreal, NIBIO), el CHM es casi plano y \texttt{peak\_local\_max} encuentra muy pocas semillas. Para mitigar este problema, se reemplaza el CHM por una \emph{superficie de densidad de copa}: para cada columna $(x, y)$ se acumula la densidad suavizada en una banda de \emph{top\_band\_m} metros bajo el v\'oxel m\'as alto de la columna. Esta superficie refleja concentraciones locales de puntos (troncos, ramas gruesas) incluso cuando la altura es uniforme, produciendo picos detectables donde el CHM no los tiene. La acumulaci\'on se implementa eficientemente mediante \emph{cumulative sum} a lo largo del eje~$Z$ del volumen suavizado, seguida de un suavizado gaussiano~2D de la superficie resultante. Sobre esta superficie se aplica \texttt{peak\_local\_max} con los mismos criterios de distancia m\'inima y umbral de altura. El par\'ametro \emph{top\_band\_m} es parte del grid search para los flujos D y E.

\paragraph{Paso~4 --- Watershed volum\'etrico 3D.} A partir de las semillas se ejecuta \texttt{skimage.segmentation.watershed} sobre el grid 3D de densidad suavizada invertida ($-\texttt{density\_smooth}$), con m\'ascara $\texttt{density\_grid} > 0$. El operador produce etiquetas de instancia por \emph{v\'oxel tridimensional}; cada punto del plot recibe la etiqueta del v\'oxel 3D que lo contiene. A diferencia de un watershed 2D sobre CHM, la asignaci\'on de puntos sigue crestas de densidad volum\'etricas y no la proyecci\'on horizontal, lo que separa correctamente troncos adyacentes y puntos de sotobosque que en un watershed 2D caer\'ian bajo la proyecci\'on de una copa vecina. Los segmentos con menos de \emph{min\_points\_per\_tree} puntos se descartan y sus puntos reciben la etiqueta 0 (fondo), que no contamina la evaluaci\'on de instancia.

\subsubsection{Calibraci\'on por grid search independiente por flujo.}
Los par\'ametros del segmentador se calibran de forma independiente para cada flujo mediante grid search exhaustivo sobre el split val, maximizando el F1 de instancia (umbral IoU 3D~$\geq 0.5$, ver Secci\'on~\ref{sec:eval}). La calibraci\'on independiente es necesaria porque los puntos de entrada al watershed difieren entre flujos: el Flujo~A recibe todos los puntos v\'alidos incluyendo suelo y vegetaci\'on baja, mientras que los dem\'as flujos reciben \'unicamente puntos clasificados como \'arbol y filtrados por HAG~$\geq 0.5$\,m. Esta diferencia en la distribuci\'on de puntos de entrada justifica un \'optimo distinto para cada flujo.

El espacio de b\'usqueda cubre las tres dimensiones que m\'as afectan la formaci\'on de semillas y la fragmentaci\'on de copas (Tabla~\ref{tab:grid-search}). Para los flujos con seeding por densidad (D y E), se a\~nade una cuarta dimensi\'on \emph{top\_band\_m} que controla el grosor de la banda superior de acumulaci\'on. Los par\'ametros restantes (\emph{min\_tree\_height} y \emph{min\_points\_per\_tree}) se mantienen fijos para mantener la dimensionalidad manejable.

\begin{table}[t]
\centering
\caption{Espacio de b\'usqueda del grid search del Watershed~3D y par\'ametros \'optimos por flujo (calibrados sobre el split val maximizando F1 de instancia con IoU 3D~$\geq 0.5$).}\label{tab:grid-search}
\resizebox{\textwidth}{!}{%
\begin{tabular}{lcccccc}
\toprule
Par\'ametro & Espacio & A & B (RF) & C (PN++) & D (RF+Dens.) & E (PN+++Dens.) \\
\midrule
voxel\_size (m)         & $\{0.1, 0.2, 0.3\}$           & 0.2 & 0.1 & 0.1 & 0.2 & 0.3 \\
gaussian\_sigma         & $\{0.3, 0.5, 1.0\}$           & 0.3 & 0.3 & 1.0 & 0.5 & 0.3 \\
min\_crown\_radius\_m   & $\{0.5, 1.0, 1.5, 2.0\}$      & 0.5 & 1.0 & 1.5 & 2.0 & 1.5 \\
top\_band\_m            & $\{2.0, 3.0, 5.0\}$           & --- & --- & --- & 2.0 & 5.0 \\
\midrule
min\_tree\_height (m)   & Fijo: 2.0                      & 2.0 & 2.0 & 2.0 & 2.0 & 2.0 \\
min\_points\_per\_tree  & Fijo: 50                       & 50  & 50  & 50  & 50  & 50  \\
\midrule
F1 val (medio)          &                                & 0.272 & 0.371 & 0.377 & 0.467 & 0.450 \\
\bottomrule
\end{tabular}%
}
\end{table}

El total de combinaciones es $3 \times 3 \times 4 = 36$ para los flujos A--C (variante CHM) y $3 \times 3 \times 4 \times 3 = 108$ para los flujos D--E (variante Density). Cada combinaci\'on se eval\'ua sobre todos los plots val, y el mejor par\'ametro por flujo es el que maximiza el F1 de instancia medio.

%% ================================================================
%%  5. PROTOCOLO DE EVALUACI\'ON
%% ================================================================
\section{Protocolo de Evaluaci\'on}\label{sec:eval}

\subsection{M\'etricas de segmentaci\'on sem\'antica}

La calidad del preprocesado sem\'antico se eval\'ua a nivel de punto individual comparando la etiqueta predicha (\'arbol / no-\'arbol) con la etiqueta ground truth derivada del treeID de FOR-instance ($\text{treeID} > 0 \rightarrow \text{\'arbol}$, excluyendo outpoints).

\textbf{tree IoU} --- Intersection over Union de la clase \'arbol:
\begin{equation}
\text{tree IoU} = \frac{TP}{TP + FP + FN}
\end{equation}

\textbf{seg mIoU} --- mean IoU promediado sobre ambas clases:
\begin{equation}
\text{seg mIoU} = \frac{1}{2}\left(\frac{TP}{TP+FP+FN} + \frac{TN}{TN+FP+FN}\right)
\end{equation}

\textbf{seg F1} --- F1 \emph{binario} de la clase positiva (\'arbol):
\begin{equation}
\text{seg F1} = \frac{2 \cdot \text{Precision} \cdot \text{Recall}}{\text{Precision} + \text{Recall}}
\end{equation}
donde $\text{Precision} = TP / (TP + FP)$ y $\text{Recall} = TP / (TP + FN)$, computados exclusivamente sobre la clase \'arbol. En la implementaci\'on esto corresponde a \texttt{sklearn.metrics.f1\_score} con \texttt{average=\textquotesingle binary\textquotesingle} y \texttt{pos\_label=1}. No se reporta macro-F1 porque el inter\'es experimental es expl\'icitamente la capacidad de cada m\'etodo para aislar la clase \'arbol, no su balance global entre clases.

\subsection{M\'etricas de segmentaci\'on de instancias}

La calidad de la segmentaci\'on de instancias se eval\'ua a nivel de \'arbol individual comparando, para cada plot, el conjunto de instancias predichas $\{P_i\}_{i=1}^{M}$ contra el conjunto de instancias ground truth $\{G_j\}_{j=1}^{N}$ derivadas del treeID de FOR-instance. Cada instancia es el subconjunto de puntos asignado a un mismo identificador. El emparejamiento sigue el protocolo establecido por ForAINet~\cite{henrich2024} y SegmentAnyTree~\cite{wielgosz2024}: a cada instancia GT se le empareja como m\'aximo una instancia predicha, eligiendo de forma \emph{greedy} la que maximiza el IoU~3D punto a punto.

\textbf{IoU 3D punto a punto.} Para cada par $(P_i, G_j)$ se define
\begin{equation}
\text{IoU}_{3D}(P_i, G_j) = \frac{|P_i \cap G_j|}{|P_i \cup G_j|},
\end{equation}
donde $|\cdot|$ denota n\'umero de puntos.

\paragraph{Matching \emph{greedy} por IoU.}
Dadas las predicciones $\{P_i\}$ y los GT $\{G_j\}$ de un plot, el procedimiento de emparejamiento es el siguiente:
\begin{enumerate}
  \item Se construye la matriz completa de IoU~3D $\mathrm{IoU}(P_i, G_j) = |P_i \cap G_j| / |P_i \cup G_j|$.
  \item Se descartan todos los pares con $\mathrm{IoU} < \tau$ con $\tau = 0{,}5$.
  \item Los pares restantes se ordenan por IoU \emph{descendente}. En caso de empate exacto se desempata por el orden natural de los identificadores (primero por menor \emph{pred\_id}, despu\'es por menor \emph{gt\_id}) para garantizar determinismo.
  \item Se recorre la lista ordenada: un par $(P_i, G_j)$ se acepta como verdadero positivo (TP) si y s\'olo si ni $P_i$ ni $G_j$ han sido asignados previamente. Tras aceptar el par, ambos se marcan como usados y no vuelven a participar en emparejamientos posteriores.
  \item Al terminar, las predicciones sin pareja cuentan como falsos positivos (FP) y los GT sin pareja como falsos negativos (FN).
\end{enumerate}
Este criterio garantiza que cada instancia predicha se asocia a lo sumo con un GT y viceversa, es determin\'istico entre corridas, y sigue el protocolo reportado en ForAINet~\cite{henrich2024} y SegmentAnyTree~\cite{wielgosz2024}. El umbral de $\tau = 0{,}5$ es el est\'andar adoptado en la literatura de segmentaci\'on de instancias forestales.

A partir de TP, FP y FN se computan precisi\'on, recall y F1 basados en el greedy matching 1-a-1 descrito arriba. Adem\'as reportamos \emph{coverage} (cobertura), definido como la fracci\'on de \'arboles GT que reciben \emph{al menos una predicci\'on} con $\text{IoU}_{3D} \geq \tau$, \textbf{sin imponer matching \'unico}. A diferencia del recall, coverage no penaliza la sobre-segmentaci\'on: si dos predicciones disjuntas solapan el mismo GT por encima del umbral, ambas contribuyen a cubrir ese GT aunque solo una pueda quedar emparejada en el greedy 1-a-1. Por construcci\'on $\text{Coverage} \geq \text{Recall}$, y la brecha $\text{Coverage} - \text{Recall}$ es un indicador directo de fragmentaci\'on de copas por el segmentador:

\begin{equation}
\text{Precision} = \frac{TP}{TP + FP}, \quad
\text{Recall} = \frac{TP}{TP + FN}
\end{equation}
\begin{equation}
\text{F1} = \frac{2 \cdot \text{Precision} \cdot \text{Recall}}{\text{Precision} + \text{Recall}}, \quad
\text{Coverage} = \frac{|\{G_j : \exists P_i \text{ con } \text{IoU}_{3D} \geq 0.5\}|}{|\{G_j\}|}
\end{equation}

Las m\'etricas se computan por plot y luego se promedian sobre los plots del split, reportando tambi\'en el conteo agregado de \'arboles GT y predichos por flujo.

\subsection{Justificaci\'on del uso de ambas tareas en cascada}

Evaluar \'unicamente la segmentaci\'on sem\'antica no permite concluir sobre la utilidad del preprocesado para la segmentaci\'on de instancias: un pipeline puede tener tree IoU alto pero F1 de instancia bajo si el segmentador no separa copas adyacentes, independientemente de la calidad de la clasificaci\'on punto a punto. Evaluar \'unicamente el F1 de instancia no permite identificar en cu\'al de las dos etapas se concentra el error. Reportar ambas tareas en cascada permite responder la pregunta central del trabajo: \emph{\'es el preprocesado sem\'antico o el segmentador el cuello de botella del pipeline?}

%% ================================================================
%%  6. RESULTADOS
%% ================================================================
\section{Resultados}\label{sec:results}

\subsection{Resultados agregados}

La Tabla~\ref{tab:results-agg} presenta los resultados agregados sobre el conjunto de test. El Flujo~D (RF + Density) alcanza el mejor F1 de instancia (0.209), superando al Flujo~B (RF + CHM, 0.142) en un 47\% relativo y al baseline (0.112) en un 87\%. El seeding por densidad mejora consistentemente ambos m\'etodos sem\'anticos: RF sube de 0.142 a 0.209 y PointNet++ de 0.082 a 0.143.

La segmentaci\'on sem\'antica de ambos clasificadores es excelente (tree IoU $>$ 0.97), pero esta calidad no se traduce proporcionalmente en calidad de instancia, evidenciando que el segmentador es el cuello de botella principal del pipeline.

\begin{table}[t]
\centering
\caption{Resultados agregados sobre el conjunto de test de FOR-instance V1. Las m\'etricas de instancia usan IoU 3D~$\geq 0.5$ como umbral de TP. Los promedios son no ponderados sobre los 11~plots.}\label{tab:results-agg}
\resizebox{\textwidth}{!}{%
\begin{tabular}{llccccccc}
\toprule
Flujo & Seeding & sem mIoU & tree IoU & sem F1 & inst Prec & inst Recall & inst F1 & Coverage \\
\midrule
A --- Sin preproc.       & CHM     & --    & --    & --    & 0.138 & 0.133 & 0.112 & 0.133 \\
B --- Random Forest      & CHM     & 0.985 & 0.990 & 0.996 & 0.394 & 0.110 & 0.142 & 0.110 \\
C --- PointNet++ MSG     & CHM     & 0.974 & 0.982 & 0.994 & 0.172 & 0.060 & 0.082 & 0.060 \\
D --- Random Forest      & Density & 0.985 & 0.990 & 0.996 & 0.418 & 0.150 & \textbf{0.209} & 0.150 \\
E --- PointNet++ MSG     & Density & 0.974 & 0.982 & 0.994 & 0.209 & 0.118 & 0.143 & 0.118 \\
\bottomrule
\end{tabular}%
}
\end{table}

\subsection{Resultados por sitio --- Segmentaci\'on sem\'antica}

La Tabla~\ref{tab:results-seg-site} muestra el tree IoU sem\'antico por sitio. Ambos clasificadores logran tree IoU $>$~0.93 en todos los sitios. RF supera a PointNet++ de forma consistente, con ventaja m\'as marcada en RMIT (0.970 vs.\ 0.940) y NIBIO (0.995 vs.\ 0.989). Los flujos D y E comparten las m\'etricas sem\'anticas de B y C respectivamente, ya que reutilizan las mismas predicciones.

\begin{table}[t]
\centering
\caption{tree IoU sem\'antico por sitio y flujo sobre el conjunto de test (promedio por sitio; no aplica al Flujo~A, que no clasifica sem\'anticamente). Los Flujos D y E comparten m\'etricas sem\'anticas con B y C.}\label{tab:results-seg-site}
\begin{tabular}{lccc}
\toprule
Sitio & GT Trees & tree IoU (B/D --- RF) & tree IoU (C/E --- PN++) \\
\midrule
CULS   & 20  & 0.995 & 0.983 \\
NIBIO  & 161 & 0.995 & 0.989 \\
RMIT   & 64  & 0.970 & 0.940 \\
SCION  & 43  & 0.998 & 0.995 \\
TUWIEN & 35  & 0.963 & 0.970 \\
\bottomrule
\end{tabular}
\end{table}

\subsection{Resultados por sitio --- Segmentaci\'on de instancias}

La Tabla~\ref{tab:results-itd-site} muestra el F1 de instancia por sitio. Los patrones m\'as relevantes son: (i)~CULS es el \'unico sitio donde el pipeline funciona razonablemente (F1 $>$ 0.5), gracias a copas bien separadas; (ii)~NIBIO (boreal denso, 50\% del test set) es el caso m\'as desafiante para todos los m\'etodos; (iii)~el seeding por densidad (Flujos D y E) mejora consistentemente en todos los sitios, con la ganancia m\'as notable en RMIT (de 0.061 a 0.213 con RF) y TUWIEN (de 0.000 a 0.136 con RF), donde el seeding CHM fallaba completamente.

\begin{table}[t]
\centering
\caption{F1 de instancia (IoU 3D~$\geq 0.5$) por sitio y flujo sobre el conjunto de test. Promedio no ponderado por sitio.}\label{tab:results-itd-site}
\resizebox{\textwidth}{!}{%
\begin{tabular}{lcccccc}
\toprule
Sitio & GT Trees & A (baseline) & B (RF+CHM) & C (PN+++CHM) & D (RF+Dens.) & E (PN+++Dens.) \\
\midrule
CULS   & 20  & 0.339 & \textbf{0.811} & 0.545 & 0.629 & 0.419 \\
NIBIO  & 161 & 0.055 & 0.090 & 0.039 & \textbf{0.151} & 0.095 \\
RMIT   & 64  & 0.241 & 0.061 & 0.000 & \textbf{0.213} & 0.194 \\
SCION  & 43  & 0.077 & 0.075 & 0.065 & \textbf{0.205} & 0.137 \\
TUWIEN & 35  & 0.168 & 0.000 & 0.000 & \textbf{0.136} & 0.115 \\
\bottomrule
\end{tabular}%
}
\end{table}

\subsection{Par\'ametros \'optimos del Watershed~3D por flujo}

Los par\'ametros \'optimos por flujo se reportan en la Tabla~\ref{tab:grid-search} (Secci\'on~\ref{sec:method}). Cabe destacar que los flujos con seeding por densidad (D y E) seleccionan voxel\_size mayores (0.2--0.3\,m) y min\_crown\_radius\_m mayores (1.5--2.0\,m) que sus contrapartes CHM, lo que sugiere que la superficie de densidad genera semillas m\'as fiables que permiten operar a mayor escala sin perder semillas cr\'iticas. El Flujo~E selecciona top\_band\_m\,=\,5.0\,m (la banda m\'as ancha), consistente con la mayor variabilidad de la sem\'antica de PointNet++ que requiere acumulaci\'on sobre un rango vertical mayor para estabilizar la superficie de densidad.

\subsection{Importancia de features --- Random Forest}

La Tabla~\ref{tab:rf-importance} muestra las cinco features m\'as importantes del clasificador RF. La HAG domina ampliamente: las dos escalas de HAG ($k\!=\!50$ y $k\!=\!20$) suman el 61.5\% de la importancia total, confirmando que la altura sobre el suelo es el predictor primario de la separaci\'on \'arbol/no-\'arbol. La verticalidad y la rugosidad a escala gruesa ($k\!=\!50$) completan el top-5, indicando que la geometr\'ia local a escala de copa aporta informaci\'on complementaria a la altura.

\begin{table}[t]
\centering
\caption{Top-5 features por importancia (Gini) en el clasificador Random Forest.}\label{tab:rf-importance}
\begin{tabular}{clc}
\toprule
Rank & Feature & Importancia \\
\midrule
1 & HAG ($k\!=\!50$)           & 0.325 \\
2 & HAG ($k\!=\!20$)           & 0.290 \\
3 & Verticalidad ($k\!=\!50$)  & 0.087 \\
4 & Rugosidad ($k\!=\!50$)     & 0.072 \\
5 & Planaridad ($k\!=\!50$)    & 0.057 \\
\bottomrule
\end{tabular}
\end{table}

\subsection{An\'alisis factorial: preprocesado $\times$ seeding}

El dise\~no factorial $3 \times 2$ permite descomponer la contribuci\'on de cada etapa. La Tabla~\ref{tab:factorial} muestra el F1 de instancia en test para cada combinaci\'on.

\begin{table}[t]
\centering
\caption{F1 de instancia (test) en el dise\~no factorial preprocesado $\times$ seeding. $\Delta$: mejora relativa del seeding Density sobre CHM.}\label{tab:factorial}
\begin{tabular}{lccc}
\toprule
Preprocesado & CHM & Density & $\Delta$ \\
\midrule
Sin preproc. (baseline) & 0.112 & ---   & ---   \\
Random Forest           & 0.142 & \textbf{0.209} & +47\% \\
PointNet++ MSG          & 0.082 & 0.143 & +74\% \\
\bottomrule
\end{tabular}
\end{table}

\paragraph{Desacoplamiento sem\'antica--instancia.}
Ambos clasificadores alcanzan tree IoU sem\'antico $>$~0.97 en todos los sitios, pero el F1 de instancia m\'aximo es 0.209 (Flujo~D). La correlaci\'on entre tree IoU y F1 de instancia es d\'ebil: RF y PointNet++ tienen tree IoU similares (0.990 vs.\ 0.982), pero su F1 de instancia difiere por un factor 2.5$\times$ con seeding CHM (0.142 vs.\ 0.082). Esto confirma que el segmentador de instancias es el cuello de botella del pipeline: mejoras en la sem\'antica por encima del 97\% de tree IoU producen retornos marginales en la calidad de instancia.

\paragraph{Efecto del seeding por densidad.}
El seeding por densidad mejora consistentemente a ambos m\'etodos sem\'anticos, con PointNet++ (+74\%) benefici\'andose m\'as que RF (+47\%) en t\'erminos relativos. Sin embargo, PointNet++ + Density (F1\,=\,0.143) apenas iguala a RF + CHM (F1\,=\,0.142), mientras que RF + Density (F1\,=\,0.209) mantiene una ventaja clara. La interacci\'on entre preprocesado y seeding es relevante: la sem\'antica m\'as limpia de RF produce menos semillas espurias, resultando en menor over-segmentaci\'on (0.200 vs.\ 0.398 para PN++ Density).

\paragraph{RF vs.\ PointNet++.}
RF supera a PointNet++ en todas las combinaciones de seeding. Esta ventaja no se explica por una sem\'antica significativamente superior (la brecha en tree IoU es $<$~1 punto porcentual), sino por una mejor interacci\'on con el segmentador: los puntos que RF clasifica err\'oneamente como \'arbol tienden a estar cerca de \'arboles reales (falsos positivos de borde de copa), mientras que los errores de PointNet++ incluyen falsos positivos aislados que generan semillas esp\'ureas.

%% ================================================================
%%  7. DISCUSI\'ON
%% ================================================================
\section{Discusi\'on}\label{sec:discussion}

\subsection{El segmentador como cuello de botella}

Los resultados muestran un patr\'on claro: la sem\'antica casi perfecta (tree IoU $>$~0.97) no se traduce en segmentaci\'on de instancias proporcional (F1 m\'aximo 0.209). El factor limitante es la capacidad del Watershed~3D para generar semillas correctas y separar copas adyacentes.

El an\'alisis factorial lo confirma: pasar de CHM a Density seeding (sin cambiar la sem\'antica) produce una mejora mayor (+47--74\%) que pasar de baseline a RF (sin cambiar el seeding, +27\%). Esto indica que la inversi\'on m\'as rentable en un pipeline ITS de dos etapas no est\'a en mejorar la sem\'antica por encima de un umbral (que RF ya supera ampliamente), sino en mejorar el segmentador de instancias.

El efecto es dependiente del sitio: en CULS (copas separadas), RF+CHM alcanza F1\,=\,0.811 y el Density seeding no aporta mejora significativa (0.629). En NIBIO (boreal denso), RMIT y TUWIEN, donde las copas se superponen, el CHM falla sistem\'aticamente (F1\,$\leq$\,0.090 con RF) y el Density seeding es necesario para obtener rendimiento no trivial. La implicaci\'on pr\'actica es que el seeding por densidad es especialmente valioso en bosques densos de altura uniforme, mientras que en bosques con copas bien diferenciadas el CHM cl\'asico es suficiente.

\subsection{Viabilidad de Random Forest sin GPU}

RF supera a PointNet++ en el pipeline completo en todas las combinaciones de seeding evaluadas: F1\,=\,0.142 vs.\ 0.082 con CHM y F1\,=\,0.209 vs.\ 0.143 con Density. Esta ventaja se obtiene sin requerir GPU: RF se entrena e infiere en CPU con scikit-learn, mientras que PointNet++ requiere entrenamiento con mixed precision. El an\'alisis de importancia de features (Tabla~\ref{tab:rf-importance}) muestra que el 61.5\% de la importancia se concentra en HAG a dos escalas, una se\~nal directamente derivable del DTM sin hardware especializado. Este resultado tiene implicaciones pr\'acticas directas: para pipelines ITS operativos donde la sem\'antica es una etapa intermedia y no el objetivo final, RF con features geom\'etricas es una alternativa m\'as eficiente y robusta que PointNet++.

\subsection{Generalizaci\'on de PointNet++ entre tipos de bosque}

PointNet++ muestra una ca\'ida de rendimiento severa en RMIT (eucaliptos, baja densidad) y TUWIEN (mixto deciduo), donde obtiene F1\,=\,0.000 con seeding CHM. Incluso con Density seeding, su F1 en estos sitios (0.194 y 0.115) es inferior al del RF (0.213 y 0.136). La menor generalizaci\'on de PointNet++ se manifiesta tambi\'en en la sem\'antica: en RMIT, PointNet++ tiene tree IoU\,=\,0.940 vs.\ 0.970 de RF, la brecha m\'as grande entre sitios. Estos resultados sugieren que PointNet++, entrenado con submuestreo global de 8{,}192 puntos, no captura adecuadamente la estructura de sitios con densidades y morfolog\'ias de copa muy diferentes a los datos de entrenamiento (dominados por NIBIO boreal). RF, al operar sobre features geom\'etricas locales sin dependencia de la distribuci\'on global, generaliza mejor entre tipos de bosque diversos.

\subsection{Limitaciones}

\begin{enumerate}
  \item \textbf{Tama\~no del conjunto de test.} El conjunto de test oficial de FOR-instance comprende 11~plots y 323~\'arboles GT. El n\'umero de \'arboles por sitio es limitado (20--161), lo que introduce varianza en las estimaciones por sitio, especialmente en CULS (20~\'arboles) y TUWIEN (35~\'arboles).

  \item \textbf{Entrenamiento \'unico de PointNet++.} Por restricciones computacionales, PointNet++ se eval\'ua con una \'unica configuraci\'on de hiperpar\'ametros. No se puede descartar que una b\'usqueda exhaustiva mejore su rendimiento.

  \item \textbf{Especificidad del segmentador.} Las conclusiones sobre el cuello de botella son espec\'ificas a las dos variantes de Watershed~3D volum\'etrico empleadas (CHM y Density seeding). Otros segmentadores de instancias forestales --- por ejemplo clustering jer\'arquico 3D, m\'etodos basados en grafos o redes neuronales end-to-end como ForAINet~\cite{henrich2024} y SegmentAnyTree~\cite{wielgosz2024} --- podr\'ian mostrar una relaci\'on diferente entre preprocesado sem\'antico y calidad de instancia.

  \item \textbf{Datos ULS de alta densidad.} Los resultados son espec\'ificos a nubes de puntos ULS con densidades en el rango $454$--$10\,156$~pts/m$^2$ (ver~\S\ref{sec:dataset}). La transferencia a datos ALS de baja densidad (${\sim}5$~pts/m$^2$) no est\'a garantizada.

  \item \textbf{Dataset de un solo benchmark.} La evaluaci\'on se realiza exclusivamente sobre FOR-instance~V1. La generalizaci\'on a otros datasets ULS o de LiDAR denso requiere validaci\'on adicional.
\end{enumerate}

\subsection{Trabajo futuro}

\begin{itemize}
  \item Comparaci\'on directa con segmentadores end-to-end (ForAINet~\cite{henrich2024}, SegmentAnyTree~\cite{wielgosz2024}) para situar el pipeline de dos etapas frente a alternativas que aprenden conjuntamente sem\'antica e instancia.
  \item Extensi\'on a FOR-instanceV2~\cite{xiang2025} para evaluar generalizaci\'on a un dataset m\'as amplio.
  \item Evaluaci\'on sobre datos ALS de baja densidad (ej.\ NeonTreeEvaluation~\cite{weinstein2021}) para determinar si el rol relativo del preprocesado sem\'antico vs.\ el segmentador cambia en reg\'imenes de baja densidad.
  \item Exploraci\'on de features radiom\'etricos adicionales (n\'umero de retorno, RGB cuando est\'e disponible) en el clasificador RF, y evaluaci\'on de su impacto sobre la segmentaci\'on sem\'antica y la calidad de instancia.
  \item Evaluaci\'on de estrategias de seeding adaptativas que combinen CHM y densidad seg\'un la variabilidad local de altura, para aprovechar las fortalezas de cada variante seg\'un el tipo de estructura forestal.
\end{itemize}

%% ================================================================
%%  8. CONCLUSIONES
%% ================================================================
\section{Conclusiones}\label{sec:conclusions}

Este trabajo evalu\'o sistem\'aticamente el impacto del preprocesado sem\'antico y la estrategia de seeding en pipelines de segmentaci\'on de \'arboles individuales (ITS) sobre datos ULS de alta densidad, mediante un dise\~no factorial $3 \times 2$ (tres preprocesados $\times$ dos variantes de Watershed~3D) sobre el benchmark FOR-instance~V1.

Los resultados muestran que: (1)~el segmentador de instancias, no el preprocesado sem\'antico, es el cuello de botella principal del pipeline --- la sem\'antica casi perfecta (tree IoU $>$~0.97) no garantiza buena segmentaci\'on de instancias; (2)~Random Forest con features geom\'etricas supera a PointNet++ MSG como preprocesado para pipelines ITS en todas las combinaciones evaluadas, sin requerir GPU; (3)~el seeding por densidad de copa mejora consistentemente al seeding cl\'asico por CHM (+47\% para RF, +74\% para PointNet++), siendo especialmente efectivo en bosques densos de altura uniforme donde el CHM es casi plano.

La combinaci\'on RF + Watershed~3D con seeding por densidad alcanza el mejor F1 de instancia (0.209 en test), estableciendo una l\'inea base robusta y eficiente para pipelines ITS de dos etapas. La brecha entre calidad sem\'antica y calidad de instancia sugiere que las mejoras m\'as impactantes en ITS no vendr\'an de clasificadores sem\'anticos m\'as sofisticados, sino de segmentadores de instancias capaces de separar copas adyacentes en canopias densas --- ya sea mediante priors sobre la geometr\'ia de copa, segmentadores end-to-end, o mejores estrategias de seeding como la propuesta en este trabajo.

%% ================================================================
%%  REFERENCIAS
%% ================================================================
\begin{thebibliography}{10}

\bibitem{puliti2023}
Puliti, S., et al.:
FOR-instance: a UAV laser scanning benchmark dataset for semantic and instance segmentation of individual trees.
arXiv preprint arXiv:2309.01279 (2023).
DOI: 10.5281/zenodo.8287792

\bibitem{qi2017}
Qi, C.R., Yi, L., Su, H., Guibas, L.J.:
PointNet++: Deep hierarchical feature learning on point sets in a metric space.
In: Advances in Neural Information Processing Systems (NeurIPS), pp.~5099--5108 (2017)

\bibitem{qi2016}
Qi, C.R., Su, H., Mo, K., Guibas, L.J.:
PointNet: Deep learning on point sets for 3D classification and segmentation.
In: IEEE Conference on Computer Vision and Pattern Recognition (CVPR), pp.~652--660 (2017)

\bibitem{yan2019pointnet2pytorch}
Yan, X.:
\texttt{Pointnet\_Pointnet2\_pytorch}: PyTorch implementation of PointNet and PointNet++.
GitHub repository, \url{https://github.com/yanx27/Pointnet_Pointnet2_pytorch} (2019). Accedido 2026-04

\bibitem{breiman2001}
Breiman, L.:
Random Forests.
Machine Learning \textbf{45}(1), 5--32 (2001)

\bibitem{weinstein2021}
Weinstein, B.G., Marconi, S., Bohlman, S., Zare, A., White, E.:
A benchmark dataset for individual tree crown delineation in co-registered airborne RGB, LiDAR and hyperspectral imagery.
Methods in Ecology and Evolution \textbf{12}(10), 1847--1856 (2021)

\bibitem{silva2016}
Silva, C.A., Hudak, A.T., Vierling, L.A., et al.:
Imputation of individual longleaf pine (\emph{Pinus palustris} Mill.) tree attributes from field and LiDAR data.
Canadian Journal of Remote Sensing \textbf{42}(5), 554--573 (2016)

\bibitem{popescu2002}
Popescu, S.C., Wynne, R.H., Nelson, R.F.:
Estimating plot-level tree heights with lidar: local filtering with a canopy-height based variable window size.
Computers and Electronics in Agriculture \textbf{37}(1--3), 71--95 (2002)

\bibitem{dalponte2016}
Dalponte, M., Coomes, D.A.:
Tree-centric mapping of forest carbon density from airborne laser scanning and hyperspectral data.
Methods in Ecology and Evolution \textbf{7}(10), 1236--1245 (2016)

\bibitem{xiang2025}
Xiang, B., et al.:
FOR-instanceV2: A large-scale benchmark for 3D forest point cloud instance segmentation.
(2025)
% [COMPLETAR: venue y detalles de publicaci\'on]

\bibitem{weinmann2015}
Weinmann, M., Jutzi, B., Hinz, S., Mallet, C.:
Semantic point cloud interpretation based on optimal neighborhoods, relevant features and efficient classifiers.
ISPRS Journal of Photogrammetry and Remote Sensing \textbf{105}, 286--304 (2015)

\bibitem{weinmann2017}
Weinmann, M., Jutzi, B., Mallet, C., Weinmann, M.:
Geometric features and their relevance for 3D point cloud classification.
ISPRS Annals of the Photogrammetry, Remote Sensing and Spatial Information Sciences \textbf{IV-1/W1}, 157--164 (2017)

\bibitem{li2023}
Li, X., et al.:
Point cloud classification in forestry using geometric features and Random Forest.
Remote Sensing \textbf{15}(8), 2103 (2023)
% [VERIFICAR cita exacta antes de submission]

\bibitem{bremer2023}
Bremer, M., et al.:
Feature importance analysis for LiDAR forest classification.
ISPRS Journal of Photogrammetry and Remote Sensing (2023)
% [VERIFICAR cita exacta antes de submission]

\bibitem{henrich2024}
Henrich, J., et al.:
ForAINet: Deep learning for forest point cloud semantic and instance segmentation.
(2024)
% [VERIFICAR venue y detalles antes de submission]

\bibitem{wielgosz2024}
Wielgosz, M., et al.:
SegmentAnyTree: A sensor and platform agnostic deep learning model for tree segmentation using laser scanning data.
(2024)
% [VERIFICAR venue y detalles antes de submission]

\bibitem{wang2023}
Wang, D., et al.:
Semantic segmentation of forest point clouds: Combining deep learning with height-normalized intensity features.
Remote Sensing \textbf{15}(8) (2023)
% [VERIFICAR cita exacta antes de submission]

% [AGREGAR: papers de ITS en ULS / LiDAR a\'ereo denso para secci\'on 2.2]
% [AGREGAR: papers de clasificaci\'on sem\'antica en ULS / dense airborne para secci\'on 2.3]
% [AGREGAR: cita para tolerancia de emparejamiento 3.0 m en ITD]
% [AGREGAR: papers de PointNet++ en datos forestales]

\end{thebibliography}

%% ================================================================
%%  AP\'ENDICE
%% ================================================================
\section*{Appendix A --- Configuraci\'on de experimentos y reproducibilidad}

\begin{table}[h]
\centering
\caption{Arquitectura detallada de PointNet++ MSG (Flujo~C). Los radios operan en el espacio normalizado del sample (tras dividir por el radio m\'aximo del bloque de 8{,}192 puntos). En cada capa SA, los canales de salida de cada grupo MSG se concatenan para formar el canal de entrada de la capa siguiente. Las capas FP usan interpolaci\'on 3-NN ponderada por $1/d$.}\label{tab:pn2-arch}
\resizebox{\textwidth}{!}{%
\begin{tabular}{llccll}
\toprule
Capa & npoint & radios & nsamples & MLPs por grupo & Salida \\
\midrule
\multicolumn{6}{l}{\emph{Encoder (Set Abstraction MSG)}} \\
SA1 & 1{,}024 & [0.1,\,0.2] & [16,\,32] & [16,16,32],\;[32,32,64] & 96\,ch \\
SA2 & 256     & [0.2,\,0.4] & [16,\,32] & [64,64,128],\;[64,96,128] & 256\,ch \\
SA3 & 64      & [0.4,\,0.8] & [16,\,32] & [128,128,256],\;[128,128,256] & 512\,ch \\
SA4 & 16      & [0.8,\,1.6] & [16,\,32] & [256,256,512],\;[256,384,512] & 1{,}024\,ch \\
\midrule
\multicolumn{6}{l}{\emph{Decoder (Feature Propagation)}} \\
FP4 & ---  & --- & --- & [512,\,512]  & 512\,ch \\
FP3 & ---  & --- & --- & [512,\,256]  & 256\,ch \\
FP2 & ---  & --- & --- & [256,\,128]  & 128\,ch \\
FP1 & ---  & --- & --- & [128,\,128]  & 128\,ch \\
\midrule
\multicolumn{6}{l}{\emph{Cabeza de clasificaci\'on}} \\
\multicolumn{6}{l}{Conv1d($128 \to 64$) + BN + ReLU + Dropout($p\!=\!0{,}5$) $\to$ Conv1d($64 \to 2$) + log\_softmax} \\
\bottomrule
\end{tabular}%
}
\end{table}

\begin{table}[h]
\centering
\caption{Detalles de reproducci\'on.}\label{tab:repro}
\begin{tabular}{ll}
\toprule
Componente & Detalle \\
\midrule
Dataset              & FOR-instance V1 (DOI: 10.5281/zenodo.8287792)  \\
Split                & Oficial dev/test sin modificaciones             \\
Hardware PointNet++  & Apple M4 Pro (memoria unificada), \texttt{batch\_size}$=16$, \texttt{bf16}, $100$ epochs; tiempo total: [pendiente medir] \\
Hardware RF          & CPU, ${\sim}16$\,min entrenamiento              \\
Framework PointNet++ & PyTorch (sin extensiones CUDA)                  \\
Framework RF         & scikit-learn                                    \\
C\'odigo             & [PENDIENTE: URL repositorio]                    \\
Semilla aleatoria    & $42$ (global: \texttt{random}, \texttt{numpy}, \texttt{torch}; ver~\S\ref{sec:splits}). La reproducibilidad bit-exact del entrenamiento de PointNet++ depende de la plataforma (MPS/CUDA/CPU) \\
\bottomrule
\end{tabular}
\end{table}

\end{document}